% \chapter{適用例}\label{cha:Application}


% 本章では、第\ref{cha:Proposal}章で作成したプロンプトを、適用した結果を示す。
% 適用対象として、「授業出席管理システム」を選定した。
% 本事例を通じて、プロンプトへの入力内容、および生成する2段階の出力（章立て構造と詳細説明）の具体像を示す。

% \section{プロンプトへの入力}

% 本プロンプトの「開発するシステム情報」に、システム概要を以下のように入力した。

% \begin{figure}[H]
% \begin{tcolorbox}[colback=white, colframe=black]
% \begin{verbatim}
% システム名: 授業出席管理システム
% 概要: 学生が自身のスマートフォンを用いて授業の出席登録を行い、教員がリアルタイムで出席状況や過去の履歴を管理・集計するシステム。
% \end{verbatim}
% \end{tcolorbox}
% \caption{適用例における開発するシステム情報}
% \label{fig:app_input}
% \end{figure}

% \section{生成結果}

% プロンプトを実行した結果、以下の出力を得た。

% \subsection{章立て（構造）}

% 出力は、章立て構造を示す。
% 画面仕様、操作フロー、機能仕様を見ると、ドメイン独自の章立て構造を生成していることが確認できる。


% \begin{figure}[H]
% \begin{tcolorbox}[
%   colback=white,
%   colframe=black,
%   fonttitle=\bfseries
% ]
% \begin{multicols}{2}
% \footnotesize
% \ttfamily
% \fontsize{11}{12}\selectfont
% \begin{verbatim}
% 1. 概要
%     1.1. 目的
%     1.2. システム概要
%     1.3. 用語定義
%     1.4. 適用範囲

% 2. システム全体像
%     2.1. 利用者区分
%     2.2. ユースケース一覧
%     2.3. システム構成図

% 3. 動作環境・基本条件
%     3.1. 動作環境
%     3.2. ネットワーク条件
%     3.3. 稼働条件
%     3.4. 制約事項

% 4. 画面仕様
%     4.1. 画面一覧
%     4.2. 学生向け画面
%         4.2.1. 出席登録画面
%         4.2.2. 出席履歴確認画面
%     4.3. 教員向け画面
%         4.3.1. 出席状況確認画面
%         4.3.2. 出席履歴管理画面

% 5. 操作フロー
%     5.1. 学生操作フロー
%     5.2. 教員操作フロー
%     5.3. 例外操作フロー

% 6. 機能仕様
%     6.1. 機能一覧
%     6.2. 出席登録機能
%     6.3. 出席状況確認機能
%     6.4. 出席履歴管理・集計機能

% 7. データ仕様
%     7.1. データ項目一覧
%     7.2. データ型・制約
%     7.3. 入出力データ定義

% 8. エラー・メッセージ仕様
%     8.1. エラー一覧
%     8.2. エラーメッセージ定義

% 9. 非機能要件
%     9.1. パフォーマンス要件
%     9.2. セキュリティ要件
%     9.3. 可用性・信頼性要件

% 10. 運用・保守
%     10.1. 通常運用
%     10.2. 障害時運用
%     10.3. データ保全

% 11. テスト観点整理
%     11.1. テスト対象範囲
%     11.2. テスト前提条件

% 12. 付帯事項
%     12.1. 前提条件
%     12.2. 制限事項
%     12.3. 将来拡張の考慮
% \end{verbatim}
% \end{multicols}
% \end{tcolorbox}
% \caption{適用例における出力}
% \label{fig:app_output1}
% \end{figure}

% \section{まとめ}

% 本適用例により、作成したプロンプトが章立てを生成できることが確認できた。